\documentclass[11pt,a4paper]{article}
\usepackage[a4paper,top=16mm,bottom=16mm,left=18mm,right=18mm]{geometry}
\usepackage[T1]{fontenc}
\usepackage{amsmath,amssymb}
\usepackage{enumitem}

\setlength{\parindent}{0pt}
\setlength{\parskip}{4pt}
\allowdisplaybreaks

\begin{document}

\begin{center}
  {\LARGE\bfseries Sequences and Equations}\\[2pt]
  {\large Answer Key}
\end{center}

\begin{enumerate}[leftmargin=*,label=\textbf{Q\arabic*.}]

\item \(n=8\). Check term: \(58\).

\item \(n=6\). Check term: \(50\).

\item
\begin{enumerate}[label=(\alph*)]
\item Heights: \(30-2n\) cm and \(21-n\) cm.
\item Same height after \(9\) hours; height \(12\) cm.
\item Candle A burns out first, after \(15\) hours.
\end{enumerate}

\item
\begin{enumerate}[label=(\alph*)]
\item Ella: \(100-3n\), Finn: \(64+6n\).
\item Same charge after \(4\) hours, at \(88\%\).
\end{enumerate}

\item
\begin{enumerate}[label=(\alph*)]
\item \(4n+1=2n+8\Rightarrow n=3.5\), not a whole number; the sequences never meet.
\item \(6n+5\) is always \(4\) more than \(6n+1\); they never meet.
\item \(c\) must be odd and at least \(3\): \(c=3,5,7,\ldots\)
\end{enumerate}

\item
\begin{enumerate}[label=(\alph*)]
\item \((2x+3)-(x+1)=(4x-1)-(2x+3)\).
\item \(x=6\); sequence \(7,15,23,\ldots\) with \(n\)th term \(8n-1\).
\end{enumerate}

\item Ali makes \(8\) houses, Bea makes \(20\) triangles, and each uses \(41\) matchsticks.

\item
\begin{enumerate}[label=(\alph*)]
\item Cost at A: \(25n+8(n+1)\) pounds.
\item \(n=10\) panels.
\item Same cost for \(2\) panels. For \(10\) panels, A is cheaper by \(\pounds16\).
\end{enumerate}

\item
\begin{enumerate}[label=(\alph*)]
\item Table: \(x,\ 3x-5,\ 9x-20,\ 27x-65\).
\item \(x=10\) cells.
\end{enumerate}

\item
\begin{enumerate}[label=(\alph*)]
\item \(n=9\) slices.
\item \(n=14.4\), not a whole number, so impossible.
\end{enumerate}

\item
\begin{enumerate}[label=(\alph*)]
\item Grey: \(n\); all slabs: \(3n+6\). Values: grey \(1,2,3,4,5,10\); all \(9,12,15,18,21,36\).
\item Fraction grey: \(\dfrac{n}{3n+6}\).
\item \(n=18\).
\item No pattern gives exactly \(\frac13\) grey; the grey fraction is always less than \(\frac13\).
\end{enumerate}

\item
\begin{enumerate}[label=(\alph*)]
\item Fraction shaded: \(\dfrac{2n-1}{2n+1}\).
\item \(n=12\).
\item \(n=19.5\), not a whole number, so impossible.
\item The fraction approaches \(1\), but never reaches \(1\); there are always \(2\) unshaded parts.
\end{enumerate}

\item Minibus: \(30\) people; coach: \(50\) people; each pays \(\pounds12\).

\item \(n\)th term: \(10n+75\). Each chair adds \(10\) cm. A \(2.05\) m stack has \(13\) chairs.

\item
\begin{enumerate}[label=(\alph*)]
\item \(3a-2b=1\) and \(a+b=7\).
\item \(a=3,\ b=4\).
\item \(3n+5:\ 8,11,14,17,20\); \quad \(4n+6:\ 10,14,18,22,26\).
\item No. Equality would give \(4(a-b)=1\), impossible for whole numbers \(a,b\).
\end{enumerate}

\item Panel costs \(\pounds30\); post costs \(\pounds12\).

\item \(x=5,\ y=7\). Sequence: \(5,7,12,19,31,50\).

\item One valid problem has \(n=7\). For example, Kai has \(n+4\) boxes of \(3\) pens, Lily has \(n\) boxes of \(5\) pens minus \(2\) pens, and they have the same number of pens. Then Lily has \(7\) boxes, Kai has \(11\) boxes, and each has \(33\) pens.

\item
\begin{enumerate}[label=(\alph*)]
\item Example: \(4n+7\) and \(97-5n\); both have 10th term \(47\).
\item No. If two different linear sequences had equal 10th and 20th terms, they would have the same common difference and same constant term, so they would be identical.
\end{enumerate}

\end{enumerate}

\end{document}